\begin{proof}[Proof of \cref{obj:thm:oracle-frontier-broad-class}]
Fix an admissible tuple \(v\). We use the constants and exponent defined in the statement.

\begin{enumerate}
\item Comparing the restrictions in
\cref{obj:def:model,obj:def:oracle-broad-model,obj:def:null,obj:def:oracle-broad-null}
gives
\[
\mathcal M_v\subseteq\mathcal M_v^{\mathrm{e,br}},
\qquad
H_0(v)\subseteq H_0^{\mathrm{e,br}}(v).
\]
The law displayed in \cref{obj:lem:causal-nonempty} therefore supplies a member of the broader model with distance \(d_0\).

For any \(P\in\mathcal M_v^{\mathrm{e,br}}\), put
\[
\bar\tau_P=\int_{[0,1]}\tau_P(x)\,\lambda(dx).
\]
Expanding the centered square and using the effect cap gives
\[
d(P)^2
=\int_{[0,1]}\tau_P(x)^2\,\lambda(dx)-\bar\tau_P^2
\le\int_{[0,1]}\tau_P(x)^2\,\lambda(dx)
\le\frac14.
\]
Thus the set of broader-model distances is nonempty and bounded above, and
\[
d_0\le D_v^{\mathrm{e,br}}\le\frac12.
\]
Furthermore, \(\gamma\le1\) implies
\[
4^{-\gamma}\ge4^{-1}=\frac14,
\qquad
c_v^{\mathrm e}
=\frac{4^{-\gamma}}{16\sqrt3}
\ge\frac1{64\sqrt3},
\qquad
C_v^{\mathrm e}>0.
\]
% lean: CausalSmith.Stat.FinitepHomogeneityDensegamma.oracle_model_distance_bounds
% lean: CausalSmith.Stat.FinitepHomogeneityDensegamma.cOracle_uniform_lower
% lean: CausalSmith.Stat.FinitepHomogeneityDensegamma.COr_pos

\item Fix an integer \(n\ge2\), and define
\[
r_{\mathrm{low}}=c_v^{\mathrm e}\rho_n^{\mathrm e}(v).
\]
By \cref{obj:lem:paired-tent-full-record-lower}, there are a null law \(P_0\in H_0(v)\) and a normalized finite prior
\[
\pi_1=\sum_{j=1}^k\omega_j\delta_{P_j},
\qquad
\omega_j\ge0,
\qquad
\sum_{j=1}^k\omega_j=1,
\]
with at least one positive weight, such that
\[
e_{P_0}(x)=\frac12,
\qquad
\omega_j>0
\ \Longrightarrow\
P_j\in\mathcal M_v,\quad
e_{P_j}(x)=\frac12,\quad
r_{\mathrm{low}}\le d(P_j)<d_0.
\]
Define the supplied function and the complete-record mixture by
\[
f=e_{P_0},
\qquad
\overline P_1^{(n)}
=\sum_{j=1}^k\omega_jP_j^{\otimes n}.
\]
The function \(f\) is continuous and satisfies \(f(x)=1/2\) everywhere. The same cited result gives
\[
\operatorname{TV}\!\left(P_0^{\otimes n},\overline P_1^{(n)}\right)<\frac12.
\]
The inclusions already established place \(P_0\) in the broader null and every positive-weight \(P_j\) in the broader model, with \(e_{P_j}=f\).

Admissibility gives \(q>0\), \(\gamma>0\), and hence \(E_0>0\). Consequently,
\[
0<r_{\mathrm{low}}
\le d(P_j)<d_0\le D_v^{\mathrm{e,br}}
\qquad(\omega_j>0).
\]
In particular, \(r_{\mathrm{low}}\) is a legal separation with a nonempty alternative class.
% lean: CausalSmith.Stat.FinitepHomogeneityDensegamma.paired_tent_full_record_lower
% lean: CausalSmith.Stat.FinitepHomogeneityDensegamma.oracle_frontier_broad_class

\item Let \(\phi^{\mathrm e}\) be an allowed supplied-propensity test satisfying the broader-null size constraint. Freeze its function input at \(f\), and average its public randomization:
\[
\overline\phi_f(\mathcal D_n)
=\int_0^1\phi^{\mathrm e}(f,\mathcal D_n,U)\,dU,
\qquad
0\le\overline\phi_f(\mathcal D_n)\le1.
\]
Joint measurability makes this a measurable function of the sample. Since all positive-weight witnesses have supplied function \(f\), integration against the finite mixture gives
\[
\int\overline\phi_f\,dP_0^{\otimes n}
=\mathsf R_n^{\mathrm e}(P_0,\phi^{\mathrm e}),
\qquad
\int\overline\phi_f\,d\overline P_1^{(n)}
=\sum_{j=1}^k\omega_j\mathsf R_n^{\mathrm e}(P_j,\phi^{\mathrm e}).
\]
For a measurable function taking values in \([0,1]\), integrating its level-set probability differences bounds the difference of its expectations by total variation. Therefore
\[
\left|
\sum_{j=1}^k\omega_j\mathsf R_n^{\mathrm e}(P_j,\phi^{\mathrm e})
-\mathsf R_n^{\mathrm e}(P_0,\phi^{\mathrm e})
\right|
\le
\operatorname{TV}\!\left(P_0^{\otimes n},\overline P_1^{(n)}\right),
\]
and hence
\[
\sum_{j=1}^k\omega_j\mathsf R_n^{\mathrm e}(P_j,\phi^{\mathrm e})
<\frac1{10}+\frac12=\frac35.
\]
If every positive-weight atom had type-II error strictly below \(2/5\), then its rejection expectation would exceed \(3/5\). Zero-weight atoms contribute zero, so normalization would imply
\[
\sum_{j=1}^k\omega_j\mathsf R_n^{\mathrm e}(P_j,\phi^{\mathrm e})
\ge\frac35\sum_{j=1}^k\omega_j
=\frac35,
\]
a contradiction. Thus
\[
\exists j:\quad
\omega_j>0,\qquad
1-\mathsf R_n^{\mathrm e}(P_j,\phi^{\mathrm e})\ge\frac25.
\]

The feasible test class is nonempty: the jointly measurable test
\[
\phi_0^{\mathrm e}(f',\mathcal D_n,U)=0
\qquad
\bigl(f'\in C([0,1]),\ \mathcal D_n\text{ an }n\text{-record sample},\
U\in[0,1]\bigr)
\]
has rejection expectation zero under every law. For every allowed test and every law,
\[
0\le\mathsf R_n^{\mathrm e}(P,\phi^{\mathrm e})\le1.
\]
Thus the alternative errors are bounded above by one. The positive-weight atom just obtained belongs to the alternative class at \(r_{\mathrm{low}}\), so
\[
\sup_{\substack{P\in\mathcal M_v^{\mathrm{e,br}}\\
d(P)\ge r_{\mathrm{low}}}}
\left\{1-\mathsf R_n^{\mathrm e}(P,\phi^{\mathrm e})\right\}
\ge\frac25.
\]
Taking the infimum over feasible tests in
\cref{obj:def:oracle-broad-risk} proves
\[
B_n^{\mathrm{e,br}}(v,r_{\mathrm{low}})\ge\frac25.
\]
This also proves the required error witness for the particular prior constructed above.
% lean: CausalSmith.Stat.FinitepHomogeneityDensegamma.oracle_single_null_prior_error_witness
% lean: CausalSmith.Stat.FinitepHomogeneityDensegamma.oracleRejectProb_bounds
% lean: CausalSmith.Stat.FinitepHomogeneityDensegamma.oracle_frontier_broad_class

\item For every separation \(r\) with \(0<r\le r_{\mathrm{low}}\), the alternative class at \(r_{\mathrm{low}}\) is contained in the alternative class at \(r\). Both are nonempty, their error suprema are bounded, and the feasible test class is unchanged. Therefore
\[
B_n^{\mathrm{e,br}}(v,r)
\ge B_n^{\mathrm{e,br}}(v,r_{\mathrm{low}})
\ge\frac25>\frac1{10}.
\]
Every successful separation in the set defining
\cref{obj:def:oracle-broad-radius} is consequently above \(r_{\mathrm{low}}\); its sentinel \(D_v^{\mathrm{e,br}}\) is also above \(r_{\mathrm{low}}\). Taking the infimum gives
\[
r_n^{*,\mathrm{e,br}}(v)\ge r_{\mathrm{low}}.
\]
% lean: CausalSmith.Stat.FinitepHomogeneityDensegamma.oracleBroadCriticalRadius_ge_of_failed

\item We now establish calibration and power directly on the broader model. Use the construction in \cref{obj:def:oracle-handle}, with
\[
s=\lfloor n/2\rfloor,\qquad
a=s^{-E_0},\qquad
T=2^{\lceil\log_2(a^{-1/(p-1)})\rceil},\qquad
J=2^{\lceil\log_2(a^{-1/\gamma})\rceil},
\]
\[
b=20T^{1-p},
\qquad
\Lambda=\frac{80T^{2-p}}s.
\]
Here
\[
s\ge1,\qquad 0<a\le1,\qquad T\ge1,\qquad J\ge1,\qquad
b\ge0,\qquad\Lambda>0.
\]
For a broader-model law \(P\), overlap makes the clipped supplied propensity equal to \(e_P\).

The elementary clipping inequalities are
\[
|y-\ell_T(y)|\le |y|^pT^{1-p},
\qquad
\ell_T(y)^2\le |y|^pT^{2-p}.
\]
Indeed, when \(|y|\le T\), the first difference vanishes and
\[
y^2=|y|^p|y|^{2-p}\le |y|^pT^{2-p}.
\]
When \(|y|>T\), the clipping remainder is at most \(|y|\), and
\[
|y|
=|y|^p|y|^{1-p}
\le |y|^pT^{1-p},
\qquad
\ell_T(y)^2\le T^2
=T^pT^{2-p}\le |y|^pT^{2-p}.
\]
The raw moment restriction gives first-moment integrability in both arms, using \(|y|\le1+|y|^p\). Integrating the two clipping inequalities yields, for \(\lambda\)-almost every \(x\),
\[
\left|\int_{\mathbb R}(y-\ell_T(y))\,Q_{\iota_{\mathrm{arm}},P}(dy\mid x)\right|
\le10T^{1-p},
\qquad
\int_{\mathbb R}\ell_T(y)^2\,Q_{\iota_{\mathrm{arm}},P}(dy\mid x)
\le10T^{2-p}
\quad(\iota_{\mathrm{arm}}\in\{0,1\}).
\]

Define the clipped conditional score mean by
\[
g_{T,P}(x)
=
\int_{\mathbb R}\ell_T(y)\,Q_{1,P}(dy\mid x)
-
\int_{\mathbb R}\ell_T(y)\,Q_{0,P}(dy\mid x).
\]
Cancellation of the arm probabilities against the inverse-propensity denominators identifies this function as the conditional mean of \(R_i\). The conditional arm-mean identities then give
\[
|g_{T,P}(x)-\tau_P(x)|\le20T^{1-p}=b
\]
almost everywhere. The conditional second moment of the score satisfies
\[
\begin{aligned}
\mathbb E(R_i^2\mid X_i=x)
&=
\frac{\int\ell_T(y)^2\,Q_{1,P}(dy\mid x)}{e_P(x)}
+
\frac{\int\ell_T(y)^2\,Q_{0,P}(dy\mid x)}{1-e_P(x)}
\\
&\le40T^{2-p}+40T^{2-p}
=80T^{2-p}.
\end{aligned}
\]
% lean: CausalSmith.Stat.FinitepHomogeneityDensegamma.oracle_clipping_moments
% lean: CausalSmith.Stat.FinitepHomogeneityDensegamma.conditional_truncation_moments

\item Define the constant-removing matrix and the sample expectation notation by
\[
u=J^{-1/2}(1,\ldots,1)^T,\qquad
C_J=\operatorname{Id}_{\mathbb R^J}-uu^T,
\]
\[
\mathbb P_{n,P}=P^{\otimes n},
\qquad
\mathbb E_{n,P}[\cdot]=\int(\cdot)\,dP^{\otimes n}.
\]
The histogram cells in \cref{obj:def:projection} are disjoint and have Lebesgue mass \(1/J\). Consequently,
\[
\int_{[0,1]}F_J(x)F_J(x)^T\,\lambda(dx)
=\operatorname{Id}_{\mathbb R^J},
\qquad
u^TF_J(x)=1,
\qquad
C_JF_J(x)=F_J(x)-u.
\]
For every \(z_{\mathrm{dir}}\in\mathbb R^J\),
\[
\|C_Jz_{\mathrm{dir}}\|^2
=\|z_{\mathrm{dir}}\|^2-(u^Tz_{\mathrm{dir}})^2
\le\|z_{\mathrm{dir}}\|^2.
\]
Combining these identities with the conditional score second-moment bound gives
\[
\begin{aligned}
\mathbb E_{n,P}
\!\left[\bigl(z_{\mathrm{dir}}^T(F_J(X_i)-u)R_i\bigr)^2\right]
&\le80T^{2-p}
\int_{[0,1]}
\bigl(z_{\mathrm{dir}}^TC_JF_J(x)\bigr)^2\,\lambda(dx)
\\
&=80T^{2-p}\|C_Jz_{\mathrm{dir}}\|^2
\\
&\le80T^{2-p}\|z_{\mathrm{dir}}\|^2.
\end{aligned}
\]

The two disjoint evaluation blocks have equal size and consist of independent records. Their common mean and centered vectors are
\[
\mu_P
=\mathbb E_{n,P}V_1
=\mathbb E_{n,P}V_2
=C_J\int_{[0,1]}F_J(x)g_{T,P}(x)\,\lambda(dx),
\qquad
\mathsf Z_{\iota_{\mathrm{block}},P}=V_{\iota_{\mathrm{block}}}-\mu_P\quad(\iota_{\mathrm{block}}\in\{1,2\}).
\]
Variance additivity within a block and the preceding single-record bound imply
\[
\mathbb E_{n,P}
\bigl[(z_{\mathrm{dir}}^T\mathsf Z_{\iota_{\mathrm{block}},P})^2\bigr]
=
\operatorname{Var}_{n,P}(z_{\mathrm{dir}}^TV_{\iota_{\mathrm{block}}})
\le\Lambda\|z_{\mathrm{dir}}\|^2.
\]
The centered vectors are independent and have mean zero.

Define the quadratic statistic by
\[
W_P=\langle V_1,V_2\rangle.
\]
Block independence gives
\[
\mathbb E_{n,P}W_P=\|\mu_P\|^2,
\qquad
W_P-\|\mu_P\|^2
=
\langle\mu_P,\mathsf Z_{1,P}\rangle
+\langle\mu_P,\mathsf Z_{2,P}\rangle
+\langle\mathsf Z_{1,P},\mathsf Z_{2,P}\rangle.
\]
All three cross-product expectations in the square of this expansion vanish:
\[
\begin{aligned}
\mathbb E_{n,P}
\bigl[\langle\mu_P,\mathsf Z_{1,P}\rangle
      \langle\mu_P,\mathsf Z_{2,P}\rangle\bigr]&=0,\\
\mathbb E_{n,P}
\bigl[\langle\mu_P,\mathsf Z_{1,P}\rangle
      \langle\mathsf Z_{1,P},\mathsf Z_{2,P}\rangle\bigr]&=0,\\
\mathbb E_{n,P}
\bigl[\langle\mu_P,\mathsf Z_{2,P}\rangle
      \langle\mathsf Z_{1,P},\mathsf Z_{2,P}\rangle\bigr]&=0.
\end{aligned}
\]
The first equality uses independence and both zero means; the other two follow by integrating first over the block appearing linearly.

For the remaining quadratic term, independence and the directional second-moment bound give
\[
\begin{aligned}
\mathbb E_{n,P}
\bigl[\langle\mathsf Z_{1,P},\mathsf Z_{2,P}\rangle^2\bigr]
&\le\Lambda\,\mathbb E_{n,P}\|\mathsf Z_{1,P}\|^2\\
&=\Lambda\sum_{j=1}^J
\mathbb E_{n,P}(\mathsf Z_{1,P})_j^2
\le J\Lambda^2.
\end{aligned}
\]
Thus
\[
\operatorname{Var}_{n,P}(W_P)
\le2\Lambda\|\mu_P\|^2+J\Lambda^2.
\]
All quantities have finite second moments, since the clipped scores and finite histogram features are bounded.

Define the good event by
\[
\mathcal G_{n,P}
=
\left\{\mathcal D_n:
|W_P-\|\mu_P\|^2|
\le16\bigl(\sqrt\Lambda\,\|\mu_P\|+\sqrt J\,\Lambda\bigr)
\right\}.
\]
Its tolerance is strictly positive. Applying the second-moment tail bound gives
\[
\begin{aligned}
\mathbb P_{n,P}(\mathcal G_{n,P}^{\,c})
&\le
\frac{2\Lambda\|\mu_P\|^2+J\Lambda^2}
{256\bigl(\sqrt\Lambda\,\|\mu_P\|+\sqrt J\,\Lambda\bigr)^2}\\
&\le\frac1{128}<0.01.
\end{aligned}
\]
The last inequality uses the nonnegative cross term in the denominator.
% lean: CausalSmith.Stat.FinitepHomogeneityDensegamma.oracleBlockVec_variance_le
% lean: CausalSmith.Stat.FinitepHomogeneityDensegamma.oracleBlockVec_inner_integral
% lean: CausalSmith.Stat.FinitepHomogeneityDensegamma.oracleBlockVec_quadratic_variance_le
% lean: CausalSmith.Stat.FinitepHomogeneityDensegamma.oracleBlockVec_tail_le

\item Suppose \(P\in H_0^{\mathrm{e,br}}(v)\), and choose its null constant:
\[
c_{\mathrm{null},P}\in[-1/2,1/2],
\qquad
\tau_P(x)=c_{\mathrm{null},P}\quad(x\in[0,1]).
\]
The histogram coefficient of this constant is \(c_{\mathrm{null},P}u\), which \(C_J\) annihilates. Hence
\[
\mu_P
=C_J\int_{[0,1]}
F_J(x)\bigl(g_{T,P}(x)-c_{\mathrm{null},P}\bigr)\,\lambda(dx).
\]
The clipping bound and cell masses imply
\[
\left|
\left(\int F_J(x)
\bigl(g_{T,P}(x)-c_{\mathrm{null},P}\bigr)\,\lambda(dx)\right)_j
\right|
\le\sqrt J\,\frac bJ=\frac b{\sqrt J}.
\]
Summing the squared coordinate bounds and using contraction by \(C_J\) gives
\[
\|\mu_P\|^2\le
\sum_{j=1}^J\frac{b^2}{J}=b^2.
\]
On \(\mathcal G_{n,P}\),
\[
W_P
\le b^2+16\sqrt\Lambda\,b+16\sqrt J\,\Lambda
\le2b^2+80\sqrt J\,\Lambda
\le2b^2+1024\sqrt J\,\Lambda.
\]
Here
\[
16\sqrt\Lambda\,b\le b^2+64\Lambda
\]
follows from \((b-8\sqrt\Lambda)^2\ge0\), and \(\sqrt J\ge1\).

The deterministic rule in \cref{obj:def:oracle-handle} rejects precisely when \(W_P\) exceeds its displayed cutoff. Its rejection event is therefore contained in \(\mathcal G_{n,P}^{\,c}\), and averaging the independent public randomization leaves that probability unchanged. Thus
\[
\mathsf R_n^{\mathrm e}(P,\phi^{*,\mathrm e}_{n,v})
\le0.01
\qquad(P\in H_0^{\mathrm{e,br}}(v)).
\]
% lean: CausalSmith.Stat.FinitepHomogeneityDensegamma.oracle_whole_null_size_le
% lean: CausalSmith.Stat.FinitepHomogeneityDensegamma.oracle_null_population_norm_le
% lean: CausalSmith.Stat.FinitepHomogeneityDensegamma.oracle_null_cutoff_le

\item For a general broader-model law \(P\), define its histogram coefficients and their associated constant by
\[
\theta_{\mathrm{hist},P}
=\int_{[0,1]}F_J(x)\tau_P(x)\,\lambda(dx),
\qquad
c_{\mathrm{hist},P}
=\frac1{\sqrt J}\sum_{j=1}^J(\theta_{\mathrm{hist},P})_j,
\qquad
\theta_{\mathrm{cen},P}=C_J\theta_{\mathrm{hist},P}.
\]
Subtracting the constant coefficient gives
\[
\theta_{\mathrm{cen},P}
=\int F_J(x)\bigl(\tau_P(x)-c_{\mathrm{hist},P}\bigr)\,\lambda(dx)
=\int F_J(x)\bigl((\Pi_J\tau_P)(x)-c_{\mathrm{hist},P}\bigr)\,\lambda(dx).
\]
The second equality holds because histogram projection preserves every cell integral.

Effect Hölder smoothness bounds the deviation from a cell average:
\[
|\tau_P(x)-(\Pi_J\tau_P)(x)|
\le20J^{-\gamma},
\qquad
\int(\tau_P-\Pi_J\tau_P)^2\,d\lambda
\le(20J^{-\gamma})^2.
\]
Indeed, every point in the same cell lies within distance \(1/J\), so averaging the Hölder bound over that cell proves the pointwise inequality.

The histogram projection is constant on each cell, and the residual has integral zero on each cell. Its residual is therefore orthogonal to the projection and to constants. Expanding the square, and using the histogram feature isometry, yields
\[
\int(\tau_P-c_{\mathrm{hist},P})^2\,d\lambda
=
\|\theta_{\mathrm{cen},P}\|^2
+\int(\tau_P-\Pi_J\tau_P)^2\,d\lambda.
\]
Also, expansion around the mean gives
\[
\int(\tau_P-c_{\mathrm{hist},P})^2\,d\lambda
=d(P)^2+(\bar\tau_P-c_{\mathrm{hist},P})^2
\ge d(P)^2.
\]
Consequently,
\[
d(P)^2
\le\|\theta_{\mathrm{cen},P}\|^2+(20J^{-\gamma})^2
\le\bigl(\|\theta_{\mathrm{cen},P}\|+20J^{-\gamma}\bigr)^2,
\]
and hence
\[
d(P)-20J^{-\gamma}\le\|\theta_{\mathrm{cen},P}\|.
\]

Finally,
\[
\mu_P-\theta_{\mathrm{cen},P}
=C_J\int F_J(x)\bigl(g_{T,P}(x)-\tau_P(x)\bigr)\,\lambda(dx).
\]
The same cellwise coefficient bound used for the null, now applied to
\(|g_{T,P}-\tau_P|\le b\), gives
\[
\|\mu_P-\theta_{\mathrm{cen},P}\|\le b.
\]
The reverse triangle inequality therefore proves the required population signal bound:
\[
d(P)-20J^{-\gamma}-b\le\|\mu_P\|.
\]
% lean: CausalSmith.Stat.FinitepHomogeneityDensegamma.centered_effect_histogram_distance
% lean: CausalSmith.Stat.FinitepHomogeneityDensegamma.oracle_effect_histogram_coefficient_distance
% lean: CausalSmith.Stat.FinitepHomogeneityDensegamma.oracle_population_clipping_error
% lean: CausalSmith.Stat.FinitepHomogeneityDensegamma.oracle_alternative_population_norm_ge

\item If
\[
d(P)\ge20J^{-\gamma}+5b+128\sqrt{\sqrt J\,\Lambda},
\]
the preceding population bound gives
\[
\|\mu_P\|\ge4b+128\sqrt{\sqrt J\,\Lambda}.
\]
Since \(J\ge1\),
\[
\sqrt\Lambda\le\sqrt{\sqrt J\,\Lambda},
\qquad
16\sqrt\Lambda\,\|\mu_P\|\le\frac{\|\mu_P\|^2}{2}.
\]
On \(\mathcal G_{n,P}\), it follows that
\[
\begin{aligned}
W_P
&\ge\|\mu_P\|^2
-16\sqrt\Lambda\,\|\mu_P\|-16\sqrt J\,\Lambda\\
&\ge\frac{\|\mu_P\|^2}{2}-16\sqrt J\,\Lambda\\
&\ge8b^2+8176\sqrt J\,\Lambda\\
&>2b^2+1024\sqrt J\,\Lambda.
\end{aligned}
\]
The penultimate inequality expands
\((4b+128\sqrt{\sqrt J\,\Lambda})^2/2\) and discards its nonnegative cross term. The last inequality is strict because \(\Lambda>0\). Thus the rule rejects throughout \(\mathcal G_{n,P}\), proving
\[
1-\mathsf R_n^{\mathrm e}(P,\phi^{*,\mathrm e}_{n,v})\le0.01
\]
at the displayed sufficient separation.
% lean: CausalSmith.Stat.FinitepHomogeneityDensegamma.oracle_power_of_mean_bound
% lean: CausalSmith.Stat.FinitepHomogeneityDensegamma.oracle_power_cutoff_lt

\item To express this separation in the public rate, define the variance-growth exponent
\[
t=\frac{2-p}{p-1}\ge0.
\]
Direct algebra gives
\[
2+t=\frac1q,
\qquad
E_0\left(2+t+\frac1{2\gamma}\right)=1.
\]
Upward dyadic rounding costs at most a factor of two:
\[
a^{-1/(p-1)}\le T\le2a^{-1/(p-1)},
\qquad
a^{-1/\gamma}\le J\le2a^{-1/\gamma}.
\]
The lower rounding bounds and the negative bias exponents imply
\[
b=20T^{1-p}\le20a,
\qquad
J^{-\gamma}\le a.
\]
Because \(0\le2-p\le1\), the upper rounding bounds imply
\[
T^{2-p}\le2a^{-t},
\qquad
\sqrt J\le\sqrt2\,a^{-1/(2\gamma)}.
\]
The exponent balance now gives
\[
s^{-1}a^{-t-1/(2\gamma)}=a^2,
\]
and consequently
\[
\sqrt J\,\Lambda
=\frac{80\sqrt J\,T^{2-p}}s
\le160\sqrt2\,a^2.
\]
Combining these estimates yields
\[
20J^{-\gamma}+5b+128\sqrt{\sqrt J\,\Lambda}
\le
\bigl(120+128\sqrt{160\sqrt2}\bigr)a.
\]

For every \(n\ge2\), the block size satisfies \(n\le3s\). Moreover,
\[
0<E_0=\frac{2\gamma q}{2\gamma+q}\le1,
\]
so
\[
a=s^{-E_0}
\le3^{E_0}n^{-E_0}
\le3\rho_n^{\mathrm e}(v).
\]
It follows that
\[
20J^{-\gamma}+5b+128\sqrt{\sqrt J\,\Lambda}
\le C_v^{\mathrm e}\rho_n^{\mathrm e}(v).
\]
Thus every broader-model law at or above the nominal upper separation has the asserted type-II bound, in particular under the stated gate
\(C_v^{\mathrm e}\rho_n^{\mathrm e}(v)<D_v^{\mathrm{e,br}}\).
All preceding estimates apply when \(s=1\) or \(J=1\); positivity of \(s,T,J,\Lambda\) suffices throughout.
% lean: CausalSmith.Stat.FinitepHomogeneityDensegamma.oracle_exponent_balance
% lean: CausalSmith.Stat.FinitepHomogeneityDensegamma.oracle_attaining_separation_le

\item Define
\[
r_{\mathrm{up}}=C_v^{\mathrm e}\rho_n^{\mathrm e}(v)>0.
\]
If \(r_{\mathrm{up}}<D_v^{\mathrm{e,br}}\), calibration makes
\(\phi^{*,\mathrm e}_{n,v}\) feasible at level \(1/10\), and its power bound gives
\[
B_n^{\mathrm{e,br}}(v,r_{\mathrm{up}})
\le
\sup_{\substack{P\in\mathcal M_v^{\mathrm{e,br}}\\d(P)\ge r_{\mathrm{up}}}}
\left\{1-\mathsf R_n^{\mathrm e}(P,\phi^{*,\mathrm e}_{n,v})\right\}
\le0.01\le\frac1{10}.
\]
For a nonempty alternative class, this supremum is bounded between zero and one; the empty-class convention assigns it zero and yields the same upper estimate. Here nonemptiness also follows from
\(r_{\mathrm{up}}<D_v^{\mathrm{e,br}}\): otherwise every model distance would be strictly below \(r_{\mathrm{up}}\), making \(r_{\mathrm{up}}\) an upper bound for their supremum. Thus \(r_{\mathrm{up}}\) belongs to the successful-radius set.

If \(D_v^{\mathrm{e,br}}\le r_{\mathrm{up}}\), the sentinel in
\cref{obj:def:oracle-broad-radius} instead gives
\[
r_n^{*,\mathrm{e,br}}(v)
\le D_v^{\mathrm{e,br}}\le r_{\mathrm{up}}.
\]
The defining set is nonempty because it contains the sentinel, and it is bounded below by zero. In both cases,
\[
r_n^{*,\mathrm{e,br}}(v)\le r_{\mathrm{up}}.
\]
Together with the lower bound, this proves the radius inequalities. The witnesses constructed above have all the required model memberships, common function input, distance bounds, normalization, total variation comparison, and positive-weight error witness.
% lean: CausalSmith.Stat.FinitepHomogeneityDensegamma.oracleBroadCriticalRadius_le_of_test
% lean: CausalSmith.Stat.FinitepHomogeneityDensegamma.oracle_frontier_broad_class

\item For \(n\ge N_v^{\mathrm e}\), the ceiling definition implies
\[
n\ge\left(\frac{2C_v^{\mathrm e}}{d_0}\right)^{1/E_0}.
\]
The base and \(E_0\) are positive. Raising this inequality to the negative power \(-E_0\) therefore gives
\[
C_v^{\mathrm e}n^{-E_0}
\le
C_v^{\mathrm e}
\left[
\left(\frac{2C_v^{\mathrm e}}{d_0}\right)^{1/E_0}
\right]^{-E_0}
=
C_v^{\mathrm e}\left(\frac{2C_v^{\mathrm e}}{d_0}\right)^{-1}
=\frac{d_0}{2}.
\]
% lean: CausalSmith.Stat.FinitepHomogeneityDensegamma.oracle_attainment_threshold

\item For \(P\in\mathcal M_v^{\mathrm{e,br}}\), the conditional arm-mean identities hold simultaneously for \(\lambda\)-almost every \(x\):
\[
\int_{\mathbb R}y\,Q_{0,P}(dy\mid x)=m_{0,P}(x),
\qquad
\int_{\mathbb R}y\,Q_{1,P}(dy\mid x)
=m_{0,P}(x)+\tau_P(x).
\]
Overlap gives \(e_P(x)>0\) and \(1-e_P(x)>0\). Pulling the constant denominators out of the integrals and cancelling them yields
\[
\begin{aligned}
&e_P(x)\int_{\mathbb R}\frac{y}{e_P(x)}\,Q_{1,P}(dy\mid x)
+(1-e_P(x))\int_{\mathbb R}\frac{-y}{1-e_P(x)}\,Q_{0,P}(dy\mid x)\\
&\qquad=
\int_{\mathbb R}y\,Q_{1,P}(dy\mid x)
-\int_{\mathbb R}y\,Q_{0,P}(dy\mid x)
=\tau_P(x).
\end{aligned}
\]
The conditional first moments are finite almost everywhere by the raw moment bound.
% lean: CausalSmith.Stat.FinitepHomogeneityDensegamma.oracle_untruncated_mean_identity

\item Finally, the coordinate maps \(v\mapsto p\) and \(v\mapsto\gamma\) are continuous in the product topology. On the admissible domain,
\[
p>0,\qquad
q=\frac{p-1}{p}>0,\qquad
2\gamma+q>0.
\]
Hence \(q\) and \(E_0\) are continuous there, and \(E_0\) is everywhere positive.

For a compact subset \(\mathcal K\) of that domain, define
\[
E_{\min,\mathcal K}
=
\begin{cases}
\displaystyle\min_{v\in\mathcal K}E_0(v),&\mathcal K\ne\varnothing,\\[1mm]
1,&\mathcal K=\varnothing.
\end{cases}
\]
In the nonempty case, continuity and compactness supply an attained minimum, whose value is positive by admissibility. In the empty case the assertions are vacuous. Thus
\[
E_{\min,\mathcal K}>0,
\qquad
E_{\min,\mathcal K}\le E_0(v)\quad(v\in\mathcal K).
\]
Together with the global lower-multiplier bound and the fixed finite upper multiplier,
\[
c_v^{\mathrm e}\ge\frac1{64\sqrt3},
\qquad
C_v^{\mathrm e}=3\bigl(120+128\sqrt{160\sqrt2}\bigr),
\]
this proves the asserted local uniformity.
% lean: CausalSmith.Stat.FinitepHomogeneityDensegamma.oracle_compact_uniformity
\end{enumerate}
\end{proof}
